\chapter{Les étapes de la démarche}
Pour chaque étape : l'objectif, le skill, les livrables, le garde-fou et le piège à éviter.

\section{Étape 0 --- Cadrage (\texttt{forge-cadrage})}
\source{skills/forge-cadrage/SKILL.md}
\textbf{Objectif.} Que chaque future fonctionnalité remonte à un objectif mesurable.\\
\textbf{Livrables.} \cmd{docs/00-cadrage/vision.md} et \cmd{impact-map.md}.

Le skill vous interroge, une question à la fois : pour qui et quel problème, situation actuelle, cadre réglementaire (APDP, OHADA, marchés publics), contraintes terrain (connectivité, terminaux, langues), puis objectifs mesurables (indicateur, valeur actuelle, cible, échéance), acteurs, changements de comportement et fonctionnalités.

L'\textbf{impact map} relie quatre niveaux : \emph{pourquoi} (objectif), \emph{qui} (acteur), \emph{comment} (changement de comportement), \emph{quoi} (fonctionnalité). Une fonctionnalité sans objectif \texttt{O\#} est signalée hors périmètre ; un chiffre inconnu est noté \emph{À mesurer}, jamais inventé. La priorisation utilise MoSCoW (Must, Should, Could, Won't).

\astuce{La validation du cadrage revient au sponsor client, pas à l'ingénieur.}

\section{Étape 1 --- Domaine (\texttt{forge-domaine})}
\source{skills/forge-domaine/SKILL.md}
\textbf{Livrables.} \cmd{event-storming.md}, \cmd{glossaire.md}, \cmd{context-map.md} dans \cmd{docs/01-domaine/}.

\begin{itemize}
  \item \textbf{Event Storming} : événements (au participe passé), commandes, acteurs, politiques, systèmes externes et points chauds. Un point chaud n'est jamais tranché sans son responsable.
  \item \textbf{Glossaire (langage omniprésent)} : un terme métier = une définition = un nom technique anglais. Les synonymes sont listés pour être \emph{bannis}. Le garde-fou \cmd{scripts/check-glossary-terms.sh} (FF-05) échoue si un synonyme banni apparaît dans le code.
  \item \textbf{Context map} : classement des sous-domaines (cœur, support, générique), regroupement en \emph{bounded contexts}, choix des patterns (couche anticorruption pour un paiement externe, Conformist pour l'IAM\dots). L'effort de conception va au cœur ; ailleurs, on achète ou on fait simple.
\end{itemize}

\section{Étape 2 --- Spécification (\texttt{forge-spec})}
\source{skills/forge-spec/SKILL.md}
\textbf{Livrables.} \cmd{docs/02-specs/SPEC-NNN.md} et \cmd{features/SPEC-NNN.feature} (Gherkin en français).

Le skill numérote la spec, copie le gabarit en \emph{Brouillon} et le remplit section par section : problème, récit utilisateur, règles métier (\emph{chacune avec sa source}), critères d'acceptation, données (données personnelles marquées), interfaces, hors périmètre, questions ouvertes.

Le fichier \cmd{.feature} contient une \cmd{Règle} par règle métier, avec au moins un cas nominal, un cas limite et un cas d'erreur.

\begin{lstlisting}
Règle: R1 — une demande complète reçoit un numéro unique
  Scénario: Soumission d'une demande complète
    Étant donné un usager authentifié
    Quand l'usager soumet la demande
    Alors un numéro au format "DEM-2026-NNNNN" lui est attribué
\end{lstlisting}

\textbf{Garde-fou.} \cmd{scripts/check-spec-validated.sh docs/02-specs/SPEC-NNN.md}. Attention : il ne contrôle \emph{que le statut}, pas le contenu. Avant de proposer \emph{Validée}, relisez que les sections 1 à 5 sont renseignées, sans marqueur \cmd{[...]} ni question ouverte bloquante, et que le \cmd{.feature} couvre chaque règle.

\section{Étape 3 --- Architecture (\texttt{forge-architecture})}
\source{skills/forge-architecture/SKILL.md}
\textbf{Livrables.} ADR dans \cmd{docs/03-architecture/adr/}, modèle C4 en Structurizr DSL (\cmd{c4/workspace.dsl}), \cmd{fitness-functions.md}.

Toute décision structurante est un ADR : \textbf{stack technique} (langage et framework back, front, mobile), base de données, hébergement, dépendance majeure, découpage. Copiez \cmd{adr/0000-template.md} en \cmd{NNNN-slug.md} ; comparez des options \emph{chiffrées} (coût en FCFA, compétences locales, exploitation par la DSI cliente). Le statut passe de \emph{Proposée} à \emph{Acceptée} par l'humain.

Le C4 est tenu \emph{en code}, jamais en image seule. Les technologies du gabarit sont des marqueurs \cmd{[STACK : à fixer par ADR]} à remplacer par celles de l'ADR de stack.

Les \textbf{fitness functions} (id \cmd{FF-NN}) empêchent l'architecture de dériver : FF-01 (imports entre modules) et FF-03 (vulnérabilités) sont branchées en CI dès le premier lot. L'ADR-0001 (monolithe modulaire) ne se contredit pas : en sortir exige un nouvel ADR qui le remplace.

\textbf{Garde-fou.} \cmd{scripts/check-adr-present.sh docs/03-architecture/adr}.

\section{Étape 4 --- Contrat d'API (\texttt{forge-contrat}) et implémentation}
\source{skills/forge-contrat, forge-implemente}
Le contrat \cmd{api/openapi.yaml} précède le code. Pour chaque opération :
\begin{itemize}
  \item \cmd{operationId} en camelCase anglais, avec les termes du glossaire ;
  \item \cmd{tags} : le bounded context ;
  \item \cmd{summary} citant \cmd{(SPEC-NNN)} --- le lien entre contrat et spécification ;
  \item erreurs en \cmd{application/problem+json} ;
  \item montants : \cmd{integer} avec suffixe \cmd{Xof} et \cmd{minimum: 0} ; dates en \cmd{date-time} (UTC).
\end{itemize}
\textbf{Garde-fou.} \cmd{scripts/check-openapi-first.sh api/openapi.yaml docs/02-specs} : chaque opération doit citer une spec existante et \emph{Validée}.

\medskip
\textbf{Implémentation (\texttt{forge-implemente}).} Avant d'écrire du code : branche correcte, spec \emph{Validée}, lecture du glossaire, des ADR acceptés et de la context map, contrat couvrant la spec. Ensuite : un scénario Gherkin = au moins un test (écrit d'abord) ; code en anglais, commentaires en français ; FCFA en entiers (jamais de \cmd{float}) ; dates en UTC au stockage et \cmd{Africa/Porto-Novo} à l'affichage ; aucun import direct entre contextes (FF-01) ; pas de dépendance majeure sans ADR. Règle absente de la spec : poser la question, ne pas deviner. Avant de conclure : linter, tests, puis \cmd{check-glossary-terms.sh}.

\section{Étape 5 --- Recette (\texttt{forge-recette})}
\source{skills/forge-recette/SKILL.md}
\textbf{Livrable.} \cmd{docs/05-recette/plan-recette.md} : un scénario Gherkin = une ligne de la table des scénarios.

\textbf{Critères d'entrée} : specs \emph{Validée}, scénarios automatisés verts en CI, jeu de données anonymisé. Sinon, la recette ne démarre pas.

\textbf{Anomalies} : bloquante (48\,h), majeure (5\,jours), mineure (lot suivant). Le classement ne change jamais sans l'avis du responsable recette client. \textbf{Sortie} : 0 bloquante, un nombre borné de majeures avec plan, exigences non fonctionnelles critiques vérifiées. Le procès-verbal est signé par le client : l'assistant ne déclare jamais la recette réussie à sa place.

Un bug trouvé en recette : reproduire, poser des hypothèses sur la cause racine, corriger le minimum via \cmd{forge-implemente}, rejouer le scénario.

\section{Étape 6 --- Production (\texttt{forge-prod})}
\source{skills/forge-prod/SKILL.md}
\textbf{Livrables.} \cmd{plan-deploiement.md}, \cmd{runbook.md}, \cmd{rollback.md}, \cmd{plan-monitoring.md}.

\textbf{Porte d'entrée} : PV de recette signé par le client et aucune anomalie bloquante. Le skill interroge ensuite sur l'hébergement réel, la capacité d'exploitation de la DSI et le budget en FCFA ; une décision d'hébergement est un ADR.

Contraintes par défaut : connexions lentes ou instables, reprise après coupure, paiements Mobile Money (file de reprise, statuts vérifiables auprès de l'opérateur), secrets en variables d'environnement. Chaque objectif \texttt{O\#} de l'impact map reçoit un indicateur, une source de mesure et une revue à J+7 et J+30.

\astuce{Le skill rédige et prépare ; il n'exécute \emph{jamais} un déploiement de production. Toute modification de CI/CD, Docker, fichiers \cmd{.env}, migration, donnée de production, \cmd{git push} ou appel à un service payant exige votre confirmation explicite.}
